\par\bigskip\pdfbookmark[0]{Über den Zusammenhang. Fragment VII}{nexus-fragment-VII}\centerline{[VII.]}\medskip
Bei einer negativen Determinante $-p$, die\begingroup\renewcommand{\thefootnote}{***)}\footnote{[d. h. wenn $p$ eine positive Primzahl der Form $4n+1$ ist.]}\endgroup\ eine Primzahl der Form $4n+1$ ist, ist die Anzahl der Klassen\begingroup\renewcommand{\thefootnote}{$\dagger$)}\footnote{[die Anzahl der Klassen ist $=2(\alpha-\beta)$.]}\endgroup\ $=(\alpha-\beta)$, wobei $\alpha$ die Anzahl der quadratischen Reste im ersten Viertel
\[
1\pdot2\pdot3\dotsfour\tfrac14(p-1)
\]
und $\beta$ die Anzahl der Nichtreste ist.
\BookRule{.30\textwidth}
